\documentclass[10pt,a4paper]{amsart}
\usepackage[margin=19mm]{geometry}
\usepackage{amsmath,amssymb,mathtools}
\usepackage[scr=rsfs,cal=euler]{mathalfa}
\usepackage{xfrac}
\usepackage[unicode,hidelinks,bookmarks=false]{hyperref}
\newcommand{\ie}{i.e.\ }
\newcommand{\cf}{cf.\ }
\newcommand{\resp}{respectively\ }
\newcommand{\Subsection}{\subsection}
\providecommand{\nobreakdash}{\nobreak\hbox{-}}
\setlength{\emergencystretch}{2em}
\multlinegap0pt
\usepackage{amssymb}
\usepackage{tikz}
\newtheorem{lem}{Lemma}
\newcommand{\cP}{{\mathcal P}}
\newcommand{\cV}{{\mathcal V}}
\title{The Geometry of Polycons and a Counterexample to Wachspress' Conjecture}
\author{Clemens Br\"user}
\date{Source: arXiv:2603.18643v1 (CC BY 4.0)}
\begin{document}
\maketitle

\noindent\emph{Source and licence.} The Geometry of Polycons and a Counterexample to Wachspress' Conjecture. Version arXiv:2603.18643v1 (CC BY 4.0), distributed by arXiv under the Creative Commons Attribution 4.0 International licence, http://creativecommons.org/licenses/by/4.0/, as recorded in the arXiv licence metadata for this exact version and shown on the abstract page https://arxiv.org/abs/2603.18643v1. Full licence notice: \texttt{licenses/arxiv-2603.18643-CC-BY-4.0.txt}.\par\smallskip

\noindent\textit{Editorial scope.} Counted unit: the article's lem resPreserv (source0.tex lines 1002--1009). The article's complete original proof of it is delivered with it (source0.tex lines 1011--1025, one \texttt{proof} environment of 15 lines). No mathematics is written, repaired or re-derived: the statement and the proof are the article's own lines, in the article's own notation, and no retained line is substituted or re-flowed.  No further source-local context is needed: every cross-reference of every retained line resolves inside the two delivered spans.  Nothing was omitted from any retained span, so the package carries no omission record.  No retained line contains a citation, so the wrapper carries no bibliography and no bibliography entry.  No retained line contains a figure environment, an image-inclusion command or drawing code, so no image is delivered and the package declares no asset dependency.  Editorial formatting only, in addition: an AMS \texttt{amsart} preamble that restates the article's own declarations and macros used by the retained lines, namely the environments \texttt{lem} on separate counters, together with the standard packages those lines need.  The retained environments print in this document as Lemma 1 in order of appearance, whereas the article prints the same items with the numbering of its own full document.  The complete native source of the article as distributed in the arXiv e-print for this exact version is bundled unchanged as \texttt{sources/196709/source0.tex}.
\begin{lem} \label{lem: resPreserv}
 Write $C_i = \cV(c_i), A_i = \cV(\alpha_i)$ and write $\cP_i$ for the polycon obtained from replacing $C_i$ with a line $L_i$. Then the following are preserved under the action of $T_\gamma$ on $\cP$:
 \begin{enumerate}
     \item $C_1, A_2$ and $A_3$.
     \item $C_1 \cap C_2$ and $C_1 \cap C_3 \cap R(\cP)$.
     \item $L_2$, i.e. the unique line through $v_{12}$ and $v_{23}$.
 \end{enumerate}
\end{lem}

\begin{proof}
 By multiplying matrices we obtain
 \begin{equation*}
     M_\cP^\gamma = T_\gamma M_{\cP}(T_\gamma)^\top = \begin{pmatrix}
         \alpha_1 + 2\gamma c_3 + \gamma^2 \alpha_2 & c_3 + \gamma \alpha_2 & c_2 + \gamma c_1 \\
         c_3 + \gamma \alpha_2 & \alpha_2 & c_1 \\
         c_2 + \gamma c_1 & c_1 & \alpha_3
     \end{pmatrix} = \begin{pmatrix}
         \alpha_1^\gamma & c_3^\gamma & c_2^\gamma \\
         c_3^\gamma & \alpha_2 & c_1 \\
         c_2^\gamma & c_1 & \alpha_3
     \end{pmatrix}.
 \end{equation*}
 From this we immediately read off the first two statements. For the third we first note that $v_{12}$ is unchanged, since $v_{12} \in C_1 \cap C_2$. Now consider the pencil of conics $c_3^\gamma = c_3 + \gamma \alpha_2$. Since $A_2$ is the adjoint of $\cP_2$, one of the pencil's base points is the unique point in $(C_3 \cap L_2) \setminus \{v_{23}\}$. Since thus two points on $L_2$ are preserved by $T_\gamma$, so is $L_2$.
\end{proof}

\end{document}
